% Options for packages loaded elsewhere
\PassOptionsToPackage{unicode}{hyperref}
\PassOptionsToPackage{hyphens}{url}
\PassOptionsToPackage{dvipsnames,svgnames,x11names}{xcolor}
\documentclass[
  10pt,
  a4paper,
]{article}
\usepackage{xcolor}
\usepackage[margin=0.8in]{geometry}
\usepackage{amsmath,amssymb}
\setcounter{secnumdepth}{-\maxdimen} % remove section numbering
\usepackage{iftex}
\ifPDFTeX
  \usepackage[T1]{fontenc}
  \usepackage[utf8]{inputenc}
  \usepackage{textcomp} % provide euro and other symbols
\else % if luatex or xetex
  \usepackage{unicode-math} % this also loads fontspec
  \defaultfontfeatures{Scale=MatchLowercase}
  \defaultfontfeatures[\rmfamily]{Ligatures=TeX,Scale=1}
\fi
\usepackage[]{newtxtext}
\ifPDFTeX\else
  % xetex/luatex font selection
\fi
% Use upquote if available, for straight quotes in verbatim environments
\IfFileExists{upquote.sty}{\usepackage{upquote}}{}
\IfFileExists{microtype.sty}{% use microtype if available
  \usepackage[]{microtype}
  \UseMicrotypeSet[protrusion]{basicmath} % disable protrusion for tt fonts
}{}
\makeatletter
\@ifundefined{KOMAClassName}{% if non-KOMA class
  \IfFileExists{parskip.sty}{%
    \usepackage{parskip}
  }{% else
    \setlength{\parindent}{0pt}
    \setlength{\parskip}{6pt plus 2pt minus 1pt}}
}{% if KOMA class
  \KOMAoptions{parskip=half}}
\makeatother
\setlength{\emergencystretch}{3em} % prevent overfull lines
\providecommand{\tightlist}{%
  \setlength{\itemsep}{0pt}\setlength{\parskip}{0pt}}

\usepackage{graphicx}
\usepackage{xurl}
\usepackage{microtype}
\usepackage{fancyhdr}
\usepackage{titlesec}
\usepackage{caption}
\usepackage{fvextra}
\usepackage{placeins}
\usepackage{needspace}
\definecolor{ink}{HTML}{173B50}
\definecolor{linkblue}{HTML}{245C75}
\captionsetup{font=small,labelfont={bf,color=ink},labelsep=period}
\titleformat{\section}{\large\bfseries\color{ink}}{}{0pt}{}
\titleformat{\subsection}{\normalsize\bfseries}{}{0pt}{}
\titlespacing*{\section}{0pt}{14pt}{6pt}
\titlespacing*{\subsection}{0pt}{10pt}{4pt}
\pagestyle{fancy}\fancyhf{}
\fancyhead[L]{\small\sffamily\color{ink}FungalLink v2}
\fancyhead[R]{\small\sffamily Preprint}
\fancyfoot[C]{\small\thepage}
\setlength{\headheight}{14pt}
\renewcommand{\headrulewidth}{0.3pt}
\setlength{\parindent}{0pt}
\setlength{\parskip}{5pt plus 1pt}
\setlength{\emergencystretch}{3em}
\setcounter{topnumber}{2}\setcounter{bottomnumber}{2}
\renewcommand{\topfraction}{0.9}\renewcommand{\textfraction}{0.08}
\renewcommand{\floatpagefraction}{0.75}
\AtBeginEnvironment{longtable}{\small}
\RecustomVerbatimEnvironment{verbatim}{Verbatim}{breaklines=true,fontsize=\footnotesize}
\AtBeginDocument{\hypersetup{colorlinks=true,linkcolor=linkblue,urlcolor=linkblue,citecolor=linkblue,breaklinks=true}}
\usepackage{bookmark}
\IfFileExists{xurl.sty}{\usepackage{xurl}}{} % add URL line breaks if available
\urlstyle{same}
% fallback for those not using the hyperref driver hyperxmp:
\makeatletter
\@ifundefined{xmpquote}{\newcommand{\xmpquote}[1]{#1}}{}
\makeatother
\hypersetup{
  colorlinks=true,
  linkcolor={Maroon},
  filecolor={Maroon},
  citecolor={Blue},
  urlcolor={linkblue},
  pdfcreator={LaTeX via pandoc}}

\author{}
\date{}

\begin{document}

\begin{center}{\LARGE\bfseries\color{ink}Supplementary Methods\par}\vspace{6pt}
{\large FungalLink v2}\\[4pt]
Vin Vomero\\[4pt]
11 October 2026\end{center}
\vspace{8pt}

Companion to FungalLink v2: a curated benchmark for linking fungal biosynthetic gene clusters to measured metabolites. References {[}1--16{]} refer to the main article. The current inventory contains 30 relationships across 16 biological query identities: 20 directly supported and ten reference-supported. Eleven relationships have complete evaluated coverage through seven query views, two have partial spectral diagnostics and seventeen lack usable inputs. Methods and results retain their original analysis versions.

\section{S1 Evidence qualification and source contexts}\label{s1-evidence-qualification-and-source-contexts}

Admission requires evidence for chemical identity at the stated resolution, an experimentally supported reference pathway, correspondence to the sampled locus or family, and a measured-target join. Direct support concerns the queried locus or an experimentally represented family member. Reference-supported assignments transfer established pathway evidence through documented correspondence. Source identification tiers remain separate from these classes. Provisional, disputed and unresolved relationships remain unlabeled. Curation includes assistant-assisted decisions under delegated authority; independent specialist re-adjudication has not been performed.

The nine GCF relationships represent five biological queries. Twenty-one BGC relationships represent eleven further queries. Multiple products of valactamide and benzomalvin retain shared query identities, while the two complete BGC-to-GCF evaluation transformations remain linked to their original biological admissions. FAC and HEx measurements do not provide native fungal co-occurrence vectors. Product-specific evidence qualifications and source discrepancies are collected in the provenance appendix.

\section{S2 Frozen task and observation populations}\label{s2-frozen-task-and-observation-populations}

The \href{https://github.com/vinvomero/fungallink-benchmark/blob/main/data/processed/source-selected-gcf-analysis-v1/manifest.json}{release manifest} defines 9,339 candidates, five queries, six products and one conservative analysis group. Candidate IDs include unnamed workbook features and remain unchanged across every workflow and sensitivity roster. The positive sets and method inputs are separate files, \href{https://github.com/vinvomero/fungallink-benchmark/blob/main/data/processed/source-selected-gcf-analysis-v1/benchmark.json}{benchmark.json} and \href{https://github.com/vinvomero/fungallink-benchmark/blob/main/data/processed/source-selected-gcf-analysis-v1/inputs.json}{inputs.json}. The source manifests, admission records, preprocessing provenance, protocols and implementation are copied under the release's \texttt{snapshots/} directory. Source provenance is indexed in \href{https://github.com/vinvomero/fungallink-benchmark/blob/main/data/processed/source-selected-gcf-analysis-v1/snapshots/data/manifests/sources.tsv}{sources.tsv}.

An external \href{https://github.com/vinvomero/fungallink-benchmark/blob/main/data/processed/source-selected-gcf-analysis-v1-results/freeze.json}{freeze receipt}, dated 7 October 2026 at 12:42:57 UTC, activates manifest SHA256 \nolinkurl{db5c9a653e873adf1c19781086ed287d65cb8836cad9a7f3bd1b6c6d33a04c84}. The receipt records verification of 123 payload files and 18 roster projections before score inspection. The immutable manifest's original \texttt{scoring\_ready:\ false} and candidate status describe its creation state; the external receipt activates that exact package for the initial analysis. Neither record constitutes a new biological admission. The package contains 89,128,545 bytes including its manifest, as recorded in the \href{https://github.com/vinvomero/fungallink-benchmark/blob/main/docs/data/source-selected-gcf-analysis-v1.md}{packaging report}.

The primary roster contains 105 source-study observations with known membership states for all five queries. These observations are not established independent strains or genomes. A genomic predictor of one denotes inclusion in the deposited source-family roster, and zero denotes covered nonmembership in that roster. Zero therefore does not establish whole-genome absence. Eleven withheld observation rows and the unallocated \texttt{NRPSLIKE\_59} member \texttt{FMGX0070} retain their unresolved status. The \href{https://github.com/vinvomero/fungallink-benchmark/blob/main/data/processed/source-selected-gcf-analysis-v1/snapshots/protocol/source-selected-gcf-analysis-decision-2026-10-07.md}{frozen operational decision} records this scope.

Historical preprocessing retains within-condition medians and recorded condition unions, maximum condition height, and detection in any required condition at a height of 10,000,000. Presence is one when any required condition is detected, zero only when all required conditions were observed below threshold, and missing otherwise. Heights remain raw; no normalization was fitted for these workflows. Missing measurements are not converted to absence.

All 18 overlapping rosters were retained: \nolinkurl{draft_all_source_observations} (105 observations), \texttt{draft\_three\_media\_only} (102), eight \nolinkurl{draft_assembly_representative_01} through \texttt{\_08} rosters (102 each), and their \texttt{\_three\_media} restrictions (99 each). The eight representative rosters enumerate all choices across three dependency pairs. They do not resolve uncertain accession relationships or prove genome independence. Exact observation IDs and projection hashes are in \href{https://github.com/vinvomero/fungallink-benchmark/blob/main/data/processed/source-selected-gcf-analysis-v1/rosters.json}{rosters.json}. Candidate pools and positives remain fixed. Sensitivity ranges describe dependence on these overlapping compositions; they are not confidence intervals or isolated causal effects of medium.

All five queries belong to \nolinkurl{caesar2023:source-selected-set-v1}, a conservative analysis group reflecting shared cohort and selection history. This grouping does not assert one biochemical pathway, and the five queries do not establish five independent pathways. In the initial six-product analysis, the directly supported stratum contains one query and two products; the reference-supported stratum contains four queries and four products. Both remain within the same analysis group.

\section{S3 Scoring workflows and permitted information}\label{s3-scoring-workflows-and-permitted-information}

The local workflows use implementation version 0.1.0 and Python's standard library. Exact definitions, input permissions, degenerate cases and code hashes are retained in the \href{https://github.com/vinvomero/fungallink-benchmark/blob/main/data/processed/source-selected-gcf-analysis-v1/snapshots/data/manifests/baseline-methods.json}{frozen method manifest}. No local workflow learns parameters or accesses evaluator positives. All scores rank higher values first, preserve exact ties and impose no score cutoff.

Let N be the number of eligible observations, g the number of source-roster members, m the number with feature detection, and o the overlap. Global prevalence is m/N. Prevalence difference is o/g − (m−o)/(N−g), with abstention if either denominator is zero. Raw Metcalf is S = 21o + N − g − 11m, equivalent to contingency weights 10 for joint presence, −10 for feature presence without membership, zero for membership without the feature, and one for joint nonmembership/nondetection. Its fixed-margin null mean is μ = 21gm/N + N − g − 11m, with variance σ² = 441gm(N−g)(N−m)/{[}N²(N−1){]}. Standardized Metcalf is (S−μ)/σ. It abstains for invariant membership; invariant features receive zero with a zero-variance flag.

These are complete-feature adaptations of Metcalf scoring, not execution of the complete NPLinker workflow. The pinned upstream reference is NPLinker v2.0.3, commit \nolinkurl{1a4603b9171b1a8e1d791cfc4b358bf51f3cfb87}; the frozen manifest records inspected files, hashes and numerical parity checks. The two Metcalf variants are related approaches. {[}2{]}

Uniform random shuffles sorted canonical candidate IDs using Python \texttt{random.Random.shuffle}. A SHA256-derived seed includes the fixed seed, repeat index, track, query type, candidate universe, cohort, candidate pool and query identities; the exact ordered fields and serialization are in the method manifest. Roster and run ID are excluded. Verified equality of those identities, candidate IDs, positive sets and analysis groups permits one shared random computation across the 18 rosters.

NPClassScore was unavailable because validated spectra and permitted chemical classes did not cover the full candidate pool. Normalized global abundance lacked normalized inputs; a phylogeny-aware control was unavailable at that analysis freeze. The subsequently completed tree and separately versioned comparison are described in S7; no pending scores are counted as completed results. Methods requiring verified genomic absence cannot use documentary roster nonmembership as a substitute. Alternative aggregation, filtered candidate universes, reconstructed-family comparisons and class/spectral sensitivities had no separately frozen inputs. These omissions describe coverage, not poor performance of unavailable methods. {[}3{]}

Source-informed family selection and clustering thresholds, known-reference use, historical rank inspection, unresolved clustering implementation/run provenance and author-workbook crosswalk provenance remain exposure limitations. Source selection by co-occurrence can favor raw Metcalf and related scores; the size of that bias was not estimated. Freezing this experiment does not erase historical exposure. No unseen-pathway or significant-superiority claim follows from the results.

\section{S4 Ranking validation and recovery metrics}\label{s4-ranking-validation-and-recovery-metrics}

A valid query output contains exactly one finite numeric score for every released candidate. Duplicate, unknown or missing candidate IDs invalidate that query output. Abstention, unavailable input, execution failure and invalid output retain distinct reasons. Required-input missingness cannot silently reduce either the candidate pool or the observation denominator. Supported product target sets are nonempty and disjoint. Multiple ions assigned to one product remain one product.

The evaluator reports any-target hit at cutoffs 1, 10 and 100, expected reciprocal rank of the first target, mean distinct-product recovery and the probability that every product is recovered. Cutoffs are capped at pool size. These measure retrieval of known positives and do not estimate biological precision or false-positive rates.

Equal numeric scores receive expected credit under uniform permutations within each tied block. For the first block containing targets, let a candidates rank above it, t candidates tie, p be targets and r = min(t, max(0, k−a)). Then \href{mailto:hit@k}{\nolinkurl{hit@k}} = 1 − C(t−p,r)/C(t,r). Expected reciprocal rank is the sum, for u = 1 through t−p+1, of C(t−u,p−1)/{[}C(t,p)(a+u){]}. It is not the reciprocal of a mean rank. Per-product hit expectations are averaged for product recovery. All-products recovery uses the joint combinatorial probability of selecting at least one candidate from every unrecovered product in the cutoff's boundary block; product probabilities are not multiplied as if independent. The \href{https://github.com/vinvomero/fungallink-benchmark/blob/main/data/processed/source-selected-gcf-analysis-v1/snapshots/scripts/evaluate_rankings.py}{frozen evaluator} and \href{https://github.com/vinvomero/fungallink-benchmark/blob/main/docs/evaluation/ranking-evaluator.md}{evaluator specification} define the counting algorithm and feasible rank bounds.

Query metrics are averaged within cohort and analysis group. Conditional estimates use valid outputs; end-to-end estimates assign zero retrieval credit to a runnable workflow's failed or unavailable query output. Whole-method unavailability receives a reason rather than fabricated scores. All four deterministic workflows covered all five queries on every roster, so conditional and end-to-end estimates coincide in this analysis.

\section{S5 Positive-set amendment and family-view extensions}\label{s5-positive-set-amendment-and-family-view-extensions}

The separate release \nolinkurl{source-selected-gcf-expanded-labels-v2} evaluates all nine admitted GCF relationships across the same five queries and 9,339 candidates. Notoamide E adds feature 4182; the pestalamide E/F-associated targets add features 1950 and 895 with their structural uncertainty retained. The notoamide query has two targets, the pestalamide query four, and the three other queries one each. Nine targets do not imply nine fully elucidated structures or independent pathways. All six BGC relationships in that dated inventory, including ilicicolin H and the three FAC targets, were excluded from this evaluation; the current ledger has 21 BGC relationships.

The \href{https://github.com/vinvomero/fungallink-benchmark/blob/main/data/processed/source-selected-gcf-expanded-labels-v2/analysis-amendment.json}{attributed amendment} followed the evidence-based admission decision and preceded inspection of the newly added targets' ranks. Original six-product results had already been viewed. This chronology supports a retrospective expansion, not a blind prospective evaluation. The new gold builder selected nine GCF relationships from the canonical ledger and preserved every initial positive, query, candidate and conservative group. The semantic track identity remains unchanged for exact cached-score reuse; the release and gold version are new and explicit.

All 72 original rankings were reused byte-for-byte with unchanged method version 0.1.0, settings and inputs. The frozen evaluator evaluated the new positive sets across four deterministic workflows and all 18 original rosters. The \href{https://github.com/vinvomero/fungallink-benchmark/blob/main/data/processed/source-selected-gcf-expanded-labels-v2-results/validation.json}{validation receipt} records full coverage, 18 verified input projections, 72 successful evaluations, zero scoring recomputations, and preservation of all 285 original package/result artifacts. Complete coverage makes conditional and end-to-end summaries identical. The directly supported stratum now contains one query with four targets; the reference-supported stratum contains four queries with five targets. All remain in one conservative analysis group.

The amendment uses exact uniform-random expectations only. No new Monte Carlo simulation was run, and the initial simulation statistics do not describe this positive set. With N candidates and p distinct singleton targets, exact \href{mailto:hit@k}{\nolinkurl{hit@k}} is 1 − C(N−k,p)/C(N,p); mean product recovery is k/N; and all-products recovery is C(k,p)/C(N,p), zero when k\textless p.~Expected reciprocal rank sums C(N−u,p−1)/{[}C(N,p)u{]} for u=1,\ldots,N−p+1. The \href{https://github.com/vinvomero/fungallink-benchmark/blob/main/data/processed/source-selected-gcf-expanded-labels-v2-results/exact-random.json}{exact-random record} preserves per-query values. Its roster invariance follows from equal candidate and positive sets, not independence between rosters.

The \href{https://github.com/vinvomero/fungallink-benchmark/blob/main/manuscript/expanded-label-results-2026-10-07.tsv}{720-row results table} contains overall, query and evidence-stratum summaries for five workflows across 18 rosters. Primary any-hit retrieval and MRR remain unchanged for the three query-dependent workflows, while all-products \href{mailto:recovery@100}{\nolinkurl{recovery@100}} drops from 1 in the initial set to 0.933333 for raw Metcalf, 0.803846 for standardized Metcalf and 0.800000 for prevalence difference. Table 2 in \href{Comparative-results.pdf}{Comparative results v1} reports the expanded primary values. A separate same-host check recomputed all 3,600 deterministic per-query metric cells with independent combinatorial code and found agreement throughout; the \href{https://github.com/vinvomero/fungallink-benchmark/blob/main/data/processed/source-selected-gcf-expanded-labels-v2-results/host-independent-verification.json}{host verification receipt} records its scope. This was automated numerical verification, not independent researcher validation. Overlapping-roster ranges remain sensitivity summaries, not biological confidence intervals; no method-superiority inference follows.

The \href{https://github.com/vinvomero/fungallink-benchmark/blob/main/data/processed/source-selected-gcf-expanded-labels-v2/evaluate_cached.py}{amendment recipe} checks identities, input projections, cached rankings and preservation before and after evaluation. Reproduction requires a restored copy and fresh results directory without its generated snapshots/manifest, because the recipe refuses to overwrite them. Gold can be rebuilt with \nolinkurl{scripts/build_source_gcf_evaluation.py} from the snapshotted ledger/group table and original inputs. The manifest pins the original evaluator and ranking dependencies; these are not newly synthesized score files. See the \href{https://github.com/vinvomero/fungallink-benchmark/blob/main/manuscript/expanded-label-results-2026-10-07.md}{expanded analysis report} for the full recipe and limits.

The \href{https://github.com/vinvomero/fungallink-benchmark/blob/main/manuscript/emericellamide-extension-results-2026-10-08.md}{separate extension report} documents the explicitly transformed HRPKS\_71 query, 19-member/86-nonmember primary roster, unchanged 9,339-candidate pool and all 18 inherited observation rosters. The underlying BGC relationship remains one biological admission. All four deterministic workflows have full candidate coverage; 72 new runs and exact single-target random expectations are retained separately from the original five-query comparisons. The frozen amendment predates these scores, while prior comparison results were already known. There is no additional independent analysis group. Supporting Figures S1--S3 preserve the original five-query comparison scope; main Figure 2 shows current readiness.

The \href{https://github.com/vinvomero/fungallink-benchmark/blob/main/manuscript/acetylaszonalenin-extension-results-2026-10-08.md}{separate extension report} records the source-numbering correction, DMAT\_34 transformation, three-member/102-nonmember primary roster and pre-score freeze. All 72 runs use the same four methods, 9,339 candidates and 18 rosters. Separate arithmetic checks cover 672,408 scores and 720 singleton metrics. Biological inventory and original comparison versions remain unchanged; evaluated coverage across separate analyses reaches 11 relationships through seven views. This does not add independent groups or establish production by every member. Reproduction commands and limitations are retained in the extension report.

\section{S6 Spectral input construction and diagnostics}\label{s6-spectral-input-construction-and-diagnostics}

The \href{https://github.com/vinvomero/fungallink-benchmark/blob/main/data/processed/trichoderma-network-candidate-v1/manifest.json}{Trichoderma candidate manifest} records one sequence-defined query, 98 candidates, two condition-level acquisitions and 91 selected spectra. Its seven spectrum-free candidates remain in the pool. The immutable candidate manifest records zero admitted labels at packaging; the later biological decision admits H without activating scoring. The ilicicolin H target is \texttt{FLTR-C000030}, historically proposed in \href{https://github.com/vinvomero/fungallink-benchmark/blob/main/data/processed/trichoderma-network-candidate-v1/proposed_labels/ilicicolin-h-proposal.tsv}{the proposal table}. That proposal remains a historical record. The later canonical admission increases the biological ledger but leaves the five scored queries unchanged.

The source is Burger et al., \href{https://doi.org/10.1186/s13068-025-02628-3}{doi:10.1186/s13068-025-02628-3}, with measurements deposited as MSV000095813. Two January 2024 TIMS-off acquisitions share a QM6a background; they cannot be counted as independent genomes. The query is the 30,604-nt QM6a reference window GL985060.1:21098--51701, not an authenticated genome of either engineered culture. Gene-model-dependent methods require a reviewed annotation policy.

The local processing used official MZmine 3.9.0 after source version 4.0.3 required login during processing. Native TSF import and local ALL-row exports were added; unsupported parameter differences are enumerated in the manifest. Source thresholds were not tuned using labels. The new candidate universe contains 98 rows after the declared filter removed 12 of 110 aligned rows. All 91 selected exported spectra were checked against native events. Native scan provenance supports computational membership. The later chemical adjudication accepts the combined source RT, fragment, competitor and sample-pattern bridge while retaining the filename-derived sample crosswalk and absence of a new standard. The \href{https://github.com/vinvomero/fungallink-benchmark/blob/main/docs/data/trichoderma-mzmine390-processing-2026-10-07.md}{processing report} preserves commands, recipes, licenses, hashes and deviations. This is an alternative-version reconstruction, not recovery of the historical network's original IDs.

\subsection{Reference-informed class diagnostics}\label{reference-informed-class-diagnostics}

Input recovery extended beyond the 11 relationships with complete evaluated coverage. A single black-yeast 6-MSA expression acquisition yielded 1,249 candidate features, of which 291 had selected MS/MS spectra. The separately retained Trichoderma ilicicolin H pool contained 98 candidates and 91 spectra. Neither pool supplied independent-genome variation for the correlation workflows used above. Experimental conditions were not treated as independent genomes.

We generated chemical-superclass inputs using the positive and negative models from the formally versioned Spec2Class software {[}16{]}, available in the \href{https://github.com/VickiPol/Spec2Class/tree/21c6f9433227c689d00549eae25e1fa2d4238fc3}{author repository}. Model files, source revisions, preprocessing and numerical environments were recorded. The positive pipeline reproduced all 100 published example final labels after recovery of a complete author fixture; the CSV version had truncated one peak array. This reproduction checks implementation compatibility, not predictive accuracy on fungal metabolites. Negative predictions covered all 291 supplied spectra. Positive predictions covered 90 of 91 supplied spectra; one had no fragments after the fixed preprocessing rules. The seven Trichoderma candidates without spectra and the empty-spectrum candidate remained uncovered. Unusual or unresolved charge annotations were retained.

A separately declared, reference-aware diagnostic combined the 6-MSA query's transferred T1PKS class with predicted superclasses using the pinned NPLinker class-association table. This broad genomic-class transfer used the known reference and was not a new antiSMASH run or an unseen-pathway evaluation. Scores were available for 291 of 1,249 candidates (23.30\%). The supported target received a score of zero: Spec2Class predicted Triterpenoids, for which the reference table contained no T1PKS association. Ninety-four candidates scored higher, and 197 candidates, including the target, shared its score. The target's rank interval was therefore 95--291 among scored candidates. Allowing the 958 unscored candidates to occupy any positions widened the full-pool interval to 95--1,249. Missing scores were not imputed, and the target result was retained without changing classes or thresholds.

A second predeclared diagnostic used the published PKS-NRPS annotation of TriliA and the retained QM6a reference correspondence for the ilicicolin H query. This literature-derived input was reference-informed, not a blind sequence prediction. The same pinned class-association kernel scored 89 of 98 candidates (90.82\%). One predicted superclass, Isoflavonoids, was absent from the reference table; together with seven missing spectra and one spectrum empty after preprocessing, this left nine candidates unscored. The supported target scored 0.00660066. Fifteen candidates scored higher and six, including the target, tied, giving a scored-only rank interval of 16--21 and conservative full-pool bounds of 16--30. These bounds describe ties and missing outputs, not statistical confidence. No unavailable score was imputed and no vocabulary remapping was used to improve the target rank.

These findings show the limits of the recovered inputs and this class-based workflow. They do not add a complete-output relationship to the evaluated count. The 6-MSA target retains its putative chemical resolution and an unresolved source-intensity discrepancy. The Trichoderma result retains reference exposure and single-background limitations. Both diagnostics remain outside the count of 11 completely evaluated relationships. The \href{https://github.com/vinvomero/fungallink-benchmark/blob/main/manuscript/tables/sixmsa-class-diagnostic-2026-10-09.tsv}{6-MSA results}, \href{https://github.com/vinvomero/fungallink-benchmark/blob/main/manuscript/tables/trichoderma-class-diagnostic-2026-10-11.tsv}{ilicicolin H results}, input protocols and run records preserve the candidate universes, exposure, missing inputs and failed attempts.

\section{S7 Marker acquisition and fixed-tree comparison}\label{s7-marker-acquisition-and-fixed-tree-comparison}

Marker acquisition finished for all 66 frozen assemblies on 8 October 2026. Inputs comprised 49 genome assemblies and 17 protein annotations. Four jobs reused pilot outputs; 62 ran within the batch. The \href{https://github.com/vinvomero/fungallink-benchmark/blob/main/manuscript/tables/marker-batch-completion-2026-10-08.tsv}{assembly-level table} records input mode, observation count, exported and eligible marker counts, completion time and retained runtime measurements. Its \href{https://github.com/vinvomero/fungallink-benchmark/blob/main/manuscript/tables/marker-batch-completion-2026-10-08.json}{provenance receipt} hashes every source quality record and completion receipt, plus each newly executed job's command and resource record.

Saved quality records contain 1,006--1,101 eligible markers per assembly (median 1,060.5). All 66 exceed the prespecified threshold of 561 of 1,122 markers. Counting accepted marker identifiers against the complete 66-assembly denominator gives 1,078 markers present in at least 53 assemblies, the ceiling of 80\% occupancy. These figures summarize acquisition records. Downstream reconstruction, alignment and site filtering subsequently completed, as described in the following fixed-tree methods. Marker counts alone do not establish a valid phylogeny; composition-test failures and fixed-model limitations remain. The 66 assemblies are observations for comparative analysis, not 66 new BGC--metabolite relationships.

Among the 62 new batch executions, 35 requested four threads and 27 requested eight. Four reused pilot jobs have no new invocation timing in this table. Runtime fields cover the final successful invocation only and exclude interrupted attempts, archival work and orchestration. Different assemblies and input modes ran under the two thread settings, so these records do not provide a controlled estimate of acceleration. The execution-cost subsection retains the scope of the earlier partial resource measurements.

The first downstream preflight stopped because its archive inventory omitted an internal BUSCO symbolic link. The corrected reader requires exact manifest membership, rejects duplicate members and invalid link chains, and retains native evidence hash checks. No marker output or biological threshold changed. The \href{https://github.com/vinvomero/fungallink-benchmark/blob/main/protocol/decisions/2026-10-08-native-archive-reader-fix.md}{reader amendment} records the failure, correction and verification. Tree construction and the version 2 comparison subsequently completed; the terminal audit and conditional interpretation are reported in the phylogeny section below.

This section records the completed tree construction and version 2 comparison. The terminal audit passed for all 7,290 distinct nonconstant pairs, including 24 recovered pairs. Complete scores and evaluation arithmetic were reconstructed from retained native logs using the original Linux interpreter.

The frozen population contains 66 observations with one distinct versioned assembly per observation. Binary genomic-family and measured-feature vectors were projected onto that population without altering cells. All five source-selected GCF queries, 9,339 measured candidates and the separate nine-target gold set were retained. Documentary family nonmembership remains distinct from demonstrated genomic absence. Comparison methods receive the same ordered vectors; labels remain confined to evaluation. The projection does not remove source selection or prior exposure.

Marker reconstruction retained 1,078 protein markers. Each marker was aligned with MAFFT 7.526 L-INS-i, using local pairwise refinement and 1,000 iterations. Missing taxa were represented by gaps. Columns required canonical amino acids in at least 50\% of retained taxa, and marker occupancy required at least 80\% of the full 66-assembly denominator. The completed concatenation contains 551,870 sites. IQ-TREE 2.4.0 inferred one unpartitioned maximum-likelihood tree with LG+F+G4, seed 20261007, four threads and a 2 GB memory limit. Partition coordinates were retained for provenance, not separately fitted marker models. No bootstrap support or other support resampling was computed.

The tree was rooted at the midpoint with Biopython 1.85. Validation required exactly 66 unique assembly tips, finite nonnegative branch lengths and a positive total length; zero branches were retained. Pairwise distances were checked across rooting and serialization. The BayesTraits adapter rescales the mean non-root branch length to 0.1. This is a comparative input tree, not a dated phylogeny or an established outgroup root. IQ-TREE reported 65 of 66 sequences failing its composition test. Fixed-model choice, midpoint rooting, compositional heterogeneity and possible gene-tree disagreement limit interpretation; none was tuned to target recovery.

The independently authored adapter invokes BayesTraits V3.0, dated 2 March 2017, with one CPU, seed 20261007, ten maximum-likelihood starts and a 120-second timeout per fit. It compares independent and dependent binary-character models. For each pair, the likelihood-ratio statistic is twice the dependent-minus-independent log likelihood, with a negative difference smaller than the 1e-6 numerical tolerance clipped to zero. A larger negative difference fails the pair. The rate quotient is \nolinkurl{(q21+q24+q31+q34)/(q12+q13+q21+q24+q31+q34+q42+q43)}. The ranking score is minus the base 10 logarithm of the chi-square survival probability with four degrees of freedom when the quotient is at least 0.65, and zero otherwise. These values implement a ranking rule; they do not support a multiple-testing-adjusted biological significance claim. Constant features receive zero, duplicate feature patterns share fits and are restored to all candidate IDs, and a constant query abstains. Invalid pair fits cause the whole query to fail without scores.

The original workflow stopped after 786 valid cached pairs when a dependent likelihood fell below its nested independent likelihood. Increasing starts to 100 did not resolve that pair. A diagnostic fixing dependent rates to the independent solution reproduced the independent likelihood at printed precision; a ten-start SBPLX diagnostic passed the existing likelihood check. A subsequent, explicitly retrospective amendment permits one SBPLX attempt for both models only after a finite, well-formed BOBYQA pair fails that check. Each model retains the higher likelihood across algorithms, with exact ties retaining BOBYQA. Selection does not use target status, rank, rate quotient or probability. Seed, starts, inputs, tree and validity thresholds remain unchanged. Other failures do not trigger this fallback; an unsuccessful fallback stops execution. This procedure does not certify a global optimum.

Version 2 writes separate outputs and reuses valid primary caches only after checking identities, trait vectors and native-file hashes. The original stopped run is preserved. The terminal audit reconstructs selected fits and complete scores from saved native logs and recomputes evaluation arithmetic without native refitting. Hash consistency and arithmetic agreement do not validate the biological labels or tree model. The final comparison contains 24 recovered pairs among 7,290 checked pairs. All five methods produced complete rankings for all five queries. The exact uniform control was evaluated separately. Full results and tie bounds are retained in \href{https://github.com/vinvomero/fungallink-benchmark/blob/main/manuscript/tables/phylogeny-comparison-v2/README.md}{the verified tables}. A Windows audit differed at the last decimal of a probability; the unchanged audit passed in the original Linux environment using byte-verified inputs. No equality threshold was relaxed.

Execution details and source locators are in the \href{https://github.com/vinvomero/fungallink-benchmark/blob/main/protocol/decisions/2026-10-07-phylogeny-comparison-population.md}{population decision}, \href{https://github.com/vinvomero/fungallink-benchmark/blob/main/protocol/decisions/2026-10-07-phylogeny-alignment-rooting.md}{alignment/rooting protocol}, \href{https://github.com/vinvomero/fungallink-benchmark/blob/main/data/manifests/phylogeny-method.json}{method manifest} and \href{https://github.com/vinvomero/fungallink-benchmark/blob/main/protocol/decisions/2026-10-11-phylogeny-optimizer-recovery.md}{recovery amendment}. The original manifest's pending-status field is historical; terminal run receipts govern execution status. The published bacterial method motivates this comparison but does not establish fungal performance {[}7{]}.

\section{S8 Similarity representations and dependence}\label{s8-similarity-representations-and-dependence}

The retained analysis environment used Python 3.13.5, RDKit 2026.3.5, Biopython 1.88, matchms 0.33.1, NumPy 2.4.6 and SciPy 1.16.3. This is distinct from the Biopython 1.85 environment used for tree rooting. Package records and hashes are retained in the companion manifest.

The frozen \href{https://github.com/vinvomero/fungallink-benchmark/blob/main/protocol/similarity-independence-v1.md}{similarity protocol} and \href{https://github.com/vinvomero/fungallink-benchmark/blob/main/data/processed/similarity-independence-v1/manifest.json}{result manifest} define the retained representations, parameters and environment. The \href{https://github.com/vinvomero/fungallink-benchmark/blob/main/docs/evaluation/similarity-independence-v1.md}{report} and \href{https://github.com/vinvomero/fungallink-benchmark/blob/main/docs/release/similarity-independence-v1.md}{dictionary} give pair-table joins and missingness rules. This companion changes no admissions, candidate universes, scores or frozen groups: the inventory remains 30 relationships/16 query identities with 11 complete, two partial and 17 missing input states.

Five supported connectivity structures yield ten numerical chemical pairs among 435 possible relationship pairs; 425 remain unavailable. RDKit Morgan fingerprints use radius 2, 2,048 bits, bond types and no chirality. No tautomer canonicalization or charge neutralization is applied. Exact isomeric equality is separate from Tanimoto equality. Source stereochemistry needs reconciliation before chiral sensitivity; independently checked source formulas are unavailable. Unresolved alternatives and proposed structures are not silently assigned exact structures.

The genomic records cover 79 source-displayed GCF members with 1,209 proteins and ten accepted anchors with 43 annotated proteins. Global Biopython alignment uses BLOSUM62, gap opening −10 and extension −0.5 with equal terminal penalties; the first library-optimal path is retained and path ties are counted. Compatible roles yield 3,575 numerical comparisons in 10,516 cross-member core pair records. Identity divides identical nongap residues by paired nongap residues; coverage divides paired residues by each full sequence length. Screens use identity at least 70\% or 90\%, each with at least 80\% coverage on both sides. Reciprocal-best ties and displayed gene order are retained. Qualifying core matches divided by all annotated proteins give lower-bound content fractions, not whole-cluster similarity. The saved reference-correspondence records retain their original evidence scope and hashes.

Spectral processing retains finite positive peaks at m/z at least 10, removes peaks within ±1 Da of known precursor m/z, and retains peaks at least 1\% of the base peak. Intensities are square-root transformed and L2-normalized; matchms CosineHungarian uses 0.02 Da tolerance, mz\_power=0 and intensity\_power=1 on the transformed intensities. All 382 retained spectra survive processing. The 4,095 Trichoderma and 42,195 black-yeast within-pool pairs total 46,290; 26,481 cross-pool pairs are excluded for opposite native-export charge signs. Cosine at least 0.7 with six matched peaks flags 93 pairs, and the 0.9 sensitivity flags 51. Native scan identity, peak-list equality and shared source-event membership remain separate fields. Source collision-energy values are 20 and 40 respectively, with units unrecorded. Missing instrument metadata remains unknown. Seven Trichoderma and 958 black-yeast candidates still have no saved native joins; no historical Caesar feature-to-native-spectrum join is supplied.

Dependence components are separately versioned scenarios built from the printed edge rules. They do not replace frozen groups or create independent observations. The conservative scenario has five components but just one with complete evaluated relationships, so an independent grouped holdout and between-group uncertainty remain unestablished. Exposure records distinguish retrospective selection, declared reference-file membership and weight-specific training membership. A pinned class table or Rosetta mapping is not a complete training roster, and no Rosetta experiment was run. Spec2Class training membership remains unknown despite retained execution receipts and weight hashes.

Software checks cover all 30 relationship IDs, eligible pair accounting, null-versus-zero handling, exact-duplicate fixtures, independent fingerprint bit arithmetic, saved alignment arithmetic and a cosine fixture of 0.48. Retained automated verification checks hashes and output structure; it does not replace human chemical or biological judgment. The v3 package carries the completed outputs and consumed inputs with an input-disposition index. No similarity calculations or retrieval experiments were rerun during manuscript integration; no new download-only replay is claimed.

\section{S9 Numerical reconstruction and execution costs}\label{s9-numerical-reconstruction-and-execution-costs}

The frozen random procedure used seed 20260930, 1,000 initial repetitions, adaptive doubling, maximum Monte Carlo standard error (MCSE) of 0.001, and a cap of 10,000 repetitions. The shared run stopped after 4,000 repetitions with maximum MCSE 0.0008148165. Exact expectations were calculated separately. MCSE describes simulation error, not biological uncertainty. One conservative analysis group provides no basis for biological bootstrap intervals or leave-one-group-out inference.

The \href{https://github.com/vinvomero/fungallink-benchmark/blob/main/data/processed/source-selected-gcf-analysis-v1-results/deterministic-summary.json}{deterministic summary} contains 72 runs: four workflows across 18 rosters. Recorded scoring and evaluation time totaled 55.60 seconds, excluding setup and serialization. The \href{https://github.com/vinvomero/fungallink-benchmark/blob/main/data/processed/source-selected-gcf-analysis-v1-results/random-resources.json}{random resource record} reports 358.55 seconds elapsed. These are execution records on the available host, not isolated method-performance benchmarks or peak-memory estimates.

A \href{https://github.com/vinvomero/fungallink-benchmark/blob/main/data/processed/source-selected-gcf-analysis-v1-results/primary-independent-audit.json}{separate numerical audit} checked all 186,780 primary scores and 80 primary any-hit/reciprocal-rank metrics without calling baseline or evaluator helpers. A \href{https://github.com/vinvomero/fungallink-benchmark/blob/main/data/processed/source-selected-gcf-analysis-v1-results/clean-deterministic-rerun.json}{clean-directory rerun} under CPython 3.13.5 with site packages disabled reproduced all 72 deterministic runs exactly as parsed JSON, including scores, statuses, audits and evaluator fields, in 63.94 seconds. This was same-host automated reproduction, not independent researcher validation; the 4,000-repeat random simulation was not rerun.

A subsequent private source rebuild used a new local Git snapshot and clean clone, cached publisher-workbook bytes and hash-pinned curation inputs. In 121.8 seconds, all 20 output files from four source stages matched exactly, the complete released input file was byte-identical, and expanded nine-target gold matched as parsed JSON. All four primary deterministic workflows reproduced their complete rankings and expanded evaluations exactly. This run used the same host and agent with isolated Python imports and pinned existing dependency files. It did not rerun the other 17 rosters or Monte Carlo control, reproduce biological adjudication, or derive the historical sample and family mappings entirely from public evidence. The \href{https://github.com/vinvomero/fungallink-benchmark/blob/main/docs/release/clean-source-reproduction-2026-10-07.md}{source-rebuild report} and \href{https://github.com/vinvomero/fungallink-benchmark/blob/main/data/processed/clean-source-reproduction-2026-10-07/host-verification.json}{host verification} define the checked stages and input boundary.

The clean-clone driver checks checkout state and all cache fingerprints before construction. It rebuilds the feature export, candidates, condition summaries and GCF inputs; builder manifests retain new runtime metadata, while the 20 scientific tables and inputs are checked exactly. A negative control changed one newline in a private copy of the sample crosswalk and was rejected before outputs were created. The first isolated-import attempt failed to find user-site openpyxl; the recorded correction copied and hashed the existing openpyxl 3.1.5 and XML-helper package files into the cache, without installing a new environment. The original failed attempt and intervention are documented. Pinned curation remains an input, including the historical FMGX-to-strain mapping; this run does not make that edge public.

The source driver and fresh-output commands are in the source-rebuild report. The following earlier recipe instead reruns from frozen final inputs; its 72-run agreement is a separate result.

To reproduce a deterministic run, use a new output directory from the repository root:

\begin{Verbatim}[breaklines=true,fontsize=\footnotesize]
$release = 'data/processed/source-selected-gcf-analysis-v1'
$rerun = 'scratch/supplement-reproduction'
New-Item -ItemType Directory -Path $rerun -ErrorAction Stop
python -S "$release/project_roster.py" --roster draft_all_source_observations --output "$rerun/inputs.json"
python -S "$release/snapshots/scripts/run_baselines.py" --inputs "$rerun/inputs.json" --method metcalf_raw --run-id supplement:primary:metcalf_raw --output "$rerun/rankings.json"
python -S "$release/snapshots/scripts/evaluate_rankings.py" --benchmark "$release/benchmark.json" --run "$rerun/rankings.json" --output "$rerun/evaluation.json"
\end{Verbatim}

Repeat over the four deterministic method IDs and 18 roster IDs without changing candidates or positives. Projection verifies input and output hashes. Source rebuilding additionally requires the pinned \texttt{openpyxl==3.1.5}; scoring and evaluation do not. Full packaging and random commands are in the \href{https://github.com/vinvomero/fungallink-benchmark/blob/main/docs/data/source-selected-gcf-analysis-v1.md}{reproduction guide}. Results and roster-specific denominators are available in the \href{https://github.com/vinvomero/fungallink-benchmark/blob/main/manuscript/initial-comparison-results-2026-10-07.tsv}{manuscript results table}.

A selected private archive reproduced all 216 deterministic method/roster combinations across the expanded five-query comparison and the two separately reported family-view extensions. Each comparison used its 18 frozen rosters and four workflows. Complete ranking and evaluation objects matched, not only headline metrics. All 458 payload hashes passed after extraction into a separate directory. Python 3.13.5 ran with isolated imports, site packages disabled and no bytecode writes; execution took 253.861 seconds. A newline change in an extracted input was rejected before scoring, and the original bytes were restored.

This extends the earlier primary-roster extension checks to all declared sensitivity rosters and packages the expanded-gold numerical stage together with them. It does not add independent replicates or new biological relationships. The replay starts from frozen numerical inputs; it does not reconstruct raw measurements or historical mappings, repeat adjudication, rerun Monte Carlo or spectral prediction, or establish independent researcher validation. It does not cover the separate completed phylogeny comparison. The \href{https://github.com/vinvomero/fungallink-benchmark/blob/main/docs/release/core-numerical-candidate-v1-2026-10-11.md}{candidate report} links the exact archive, receipts and commands. Those exact numerical package bytes are included in the public version 1.0.0 release with the scoped source notice; historical candidate metadata is preserved.

The \href{https://github.com/vinvomero/fungallink-benchmark/blob/main/manuscript/resource-costs-2026-10-07.tsv}{resource-cost table} and \href{https://github.com/vinvomero/fungallink-benchmark/blob/main/manuscript/resource-costs-2026-10-07.md}{measurement report} record retained timing, memory and storage receipts. Original scoring and evaluation across 72 method/roster combinations used 55.603 seconds summed over the measured intervals and 55.188 CPU-seconds. A later exact 72-run reproducibility check took 63.938 wall seconds. The initial shared 4,000-repeat random computation took 358.546 wall seconds and 355.297 CPU-seconds; its result was reused across 18 rosters and reached the requested precision (maximum Monte Carlo standard error 0.0008148165324141222, tolerance 0.001; receipt status \texttt{precision\_achieved}). The 4,000 repeats were below the configured 10,000-repeat maximum. Expanded-gold evaluation reused all 72 rankings and took 16.696 wall seconds. Clean-source reproduction of its declared stages took 121.795 wall seconds, with 216,760,145 bytes across 38 receipt-listed output files. These runs do not measure new biological curation or an independent public-only reconstruction.

Completed real-data marker receipts cover five protein pilots, one genome pilot and three new batch jobs. Protein inference took 107.98--149.47 seconds per completed job; the three genome jobs took 1,109.53--1,528.08 seconds. Each requested four CPUs. The genome pilot reused by the batch is counted only once. Its GNU time measurement was 1,109.53 wall seconds, 4,388.00 CPU-seconds and 979,329,024 bytes maximum RSS. Incomplete jobs are excluded. The six-taxon synthetic native-tree smoke provides only a component timing of 5.783 wall seconds for IQ-TREE; it does not estimate fungal tree runtime.

Wall time and CPU time remain separate. Linux maximum RSS is converted from GNU time kbytes using 1,024 bytes per unit; Windows peak working set is a cumulative process measurement, not an isolated method peak. Logical output and compressed archive sizes are separate disk measures. Missing values are \texttt{NA}, including temporary peak storage and clean-source CPU/RAM. This is a measured computational subset, not total benchmark, curator, full-phylogeny or monetary cost. The linked report documents completion evidence, versions, hardware limits and exclusions without estimating unmeasured stages. This historical subset predates the completed phylogeny comparison described in S7.

\section{S10 Record directory}\label{s10-record-directory}

Paths below identify released records or retained historical artifacts. The public release README, manifest and evidence crosswalk distinguish bundled files, pinned upstream retrieval and omitted source materials. Use the full data archive for numerical reconstruction.

\subsection{Evidence ledger}\label{evidence-ledger}

\nolinkurl{data/curated/label-evidence.tsv}

\subsection{Readiness table}\label{readiness-table}

\nolinkurl{data/curated/relationship-readiness-v1.tsv}

\subsection{Readiness dictionary}\label{readiness-dictionary}

\nolinkurl{docs/release/relationship-readiness-v1.md}

\subsection{Task and exposure manifest}\label{task-and-exposure-manifest}

\nolinkurl{data/manifests/task-exposure-v1.tsv}

\subsection{Evidence and payload manifest}\label{evidence-and-payload-manifest}

\nolinkurl{data/manifests/release-evidence-v1.tsv}

\subsection{Release map}\label{release-map}

\nolinkurl{docs/release/release-map-v1.md}

\subsection{Reproduction guide}\label{reproduction-guide}

\nolinkurl{docs/release/reproduction-stages-v1.md}

\subsection{Reproduction stages table}\label{reproduction-stages-table}

\nolinkurl{docs/release/reproduction-stages-v1.tsv}

\subsection{Similarity output directory}\label{similarity-output-directory}

\nolinkurl{data/processed/similarity-independence-v1/}

\subsection{Similarity dictionary}\label{similarity-dictionary}

\nolinkurl{docs/release/similarity-independence-v1.md}

\subsection{Scientific review worksheet}\label{scientific-review-worksheet}

\nolinkurl{manuscript/scientific-review-v1.tsv}

\section{Appendix A Evidence and analytical provenance}\label{appendix-a-evidence-and-analytical-provenance}

This retained chronology documents the source and admission boundaries used by the analyses. The current readiness view takes precedence over intermediate counts.

Admission records distinguish chemical identity, experimental reference-pathway support, sampled-locus or family correspondence, and measured-target correspondence. Direct support requires an experimentally represented locus or family member joined to the measured product. Reference-supported assignments transfer established pathway evidence through justified correspondence. Provisional and unresolved links remain unlabeled. Source chemical-identification tiers and spectral availability do not replace these biological criteria.

The current biological inventory contains 30 admitted relationships across 16 query identities: nine GCF-product relationships across five GCF queries and 21 BGC-product relationships across eleven BGC queries. Twenty relationships are directly supported and ten are reference-supported. On 8 October, seven valactamide B--H-associated measured products were admitted after review of their product-specific deletion comparisons (Supplementary Figures 20--22 of Clevenger et al.~{[}10{]}). Their structures remain proposed, and all share the already admitted valactamide query. A printed mass discrepancy for B is retained; its named historical panel is not an automatic raw-feature join. The \href{https://github.com/vinvomero/fungallink-benchmark/blob/main/protocol/decisions/2026-10-08-valactamide-congener-admissions.md}{dated decision} records the evidence scope. Evaluated coverage remains 11 relationships across the separate analyses; all fifteen FAC relationships lack complete evaluated coverage.

Two subsequent FAC admissions add distinct queries: reference-supported acu-dioxomorpholine A {[}11{]} and directly supported terreazepine {[}12{]}. Acu-dioxomorpholine uses a source-identified reference interval, with the tested-clone correspondence limit and gene-direction discrepancy retained. Terreazepine uses the deposited tested FAC MT376588.1; its reported S/R mixture is one measured target. Both retain documented source/deposit discrepancies and incomplete candidate-pool and sample/feature joins. Neither adds evaluated coverage.

The HEx study adds two reference-supported relationships: PKS17 to source-characterized compound 6 and PKS15 to the source-identified orsellinic-acid-associated compound 3 {[}13{]}. All 13 and seven annotated reference proteins, respectively, match the archived expression records. The measurements come from engineered yeast, so they are kept in a separate expression cohort and do not supply native fungal co-occurrence vectors. Complete candidate tables and native feature joins remain unresolved. The two HEx admissions add two queries without changing evaluated coverage.

One further directly supported relationship is the free linear Anth-NmPhe-Anth intermediate measured from intact AtFAC9J20 and compared with a synthetic standard by tandem mass spectrometry {[}14{]}. It shares the benzomalvin query and adds neither an independent pathway nor evaluated coverage. Its source-panel identity is retained until native feature joins are available.

The black-yeast expression study adds one directly supported relationship between a deposited fungal synthase expression CDS and the 6-MSA-associated measured product {[}15{]}. The queried locus is in the sampled S. cerevisiae expression strain, whose BioSample alias links its assembly to the source acquisition. Its translated sequence differs by one amino acid from the cited reference; the observed sequence is retained. The chemical name remains putative, and the label denotes the measured source panel rather than exact-structure gold. This adds one query in a separate expression cohort, whose recovered full candidate pool is assessed separately in the spectral diagnostic in S6.

Two further benzomalvin-pathway measured targets, the macrocycle-associated compound 2 and desmethyl-associated compound 4 from Clevenger et al.~2018 {[}14{]}, were admitted at putative chemical resolution. Both share the existing tested FAC query and add no query identity, pathway group or evaluated case. The \href{https://github.com/vinvomero/fungallink-benchmark/blob/main/protocol/decisions/2026-10-09-benzomalvin-two-measured-targets.md}{admission decision} preserves the product-specific evidence and source-panel limitations. These counts describe the present evidence inventory, not a final benchmark cohort or the scored population.

On 7 October 2026, the research lead delegated the remaining admission judgments to the assistant. Seven relationships were then admitted through assistant-curated evidence judgments for final manuscript review; this was not independent human or external scientific review.

The frozen initial experiment predates those seven admissions and retains five GCF queries and six products. Its directly supported query comprises pestalamides A and B; its four reference-supported queries concern aspernidine A, fellutamide B, griseofulvin and notoamide B. The three additional GCF labels, notoamide E and the structurally putative pestalamide E- and F-associated features, are evaluated in the separate nine-target amendment in S5. They do not retroactively change the initial positive sets or scores.

The four BGC relationships admitted on 7 October are ilicicolin H, supported by Burger et al.~{[}9{]}, and three historical FAC targets from Clevenger et al.~{[}10{]}: benzomalvin A/D, the unnamed measured sesterterpenoid facms0005, and valactamide A. Ilicicolin H has a separate candidate package containing 98 features and 91 spectra, with a separately declared partial class-linking diagnostic in S6. The three FAC relationships have incomplete benchmark inputs; their admission does not validate joins to the native-donor pilot. Biological eligibility, method-input readiness and analytical activation remain separate. The dated \href{https://github.com/vinvomero/fungallink-benchmark/blob/main/protocol/decisions/2026-10-07-delegated-admission-dispositions.md}{admission decision} and Supplementary Methods record the scopes and limitations.

\end{document}
